%% ============================================================
\section{Test Harness}\label{sec:harness}
%% ============================================================

Executing tests on real hardware is not a trivial endeavor, as we must run the test case in supervisor mode while limiting potential interference from other programs and the rest of the operating system. 
Further, we must carefully eliminate interference from the test case itself; e.g., we cannot store the result of a memory access, as this adds additional memory accesses that might interfere with test execution.

Instead of writing our own test harness from scratch, we leverage Barrelfish's multikernel architecture~\cite{Baumann:2009:Multikernel}. Specifically, it has isolated per-core kernels and runs system calls with interrupts disabled, limiting interference.
That lets us leverage Barrelfish's infrastructure for booting cores and memory management.
We integrate the test cases into the kernel. We write a small user-space driver that performs setup, coordinates among the cores, and then invokes the test case in the kernel via a system call.
At the end, we collect the results: the observed memory accesses in the test program (i.e., whether they caused a fault and the value read).

We cannot use shared-memory barriers for synchronization among cores, as they introduce memory accesses that add undesired interference with the tests. Instead, we leverage synchronized time-stamp counters for synchronization. The test drivers agree on which timestamp counter value each core starts its execution.
